\documentclass[12pt, a4paper]{book}

% ─── Core packages ────────────────────────────────────────────────────────────
\usepackage[T1]{fontenc}
\usepackage[utf8]{inputenc}
\usepackage[charter, cal=cmcal]{mathdesign}
\usepackage[scaled=0.92]{helvet}
\usepackage{geometry}
\geometry{
  a4paper,
  top=30mm, bottom=30mm,
  inner=32mm, outer=25mm,
  marginparwidth=18mm
}

% ─── Mathematics ──────────────────────────────────────────────────────────────
\usepackage{amsmath, amsthm}

% ─── Chemistry ────────────────────────────────────────────────────────────────
\usepackage[version=4]{mhchem}

% ─── Graphics & colour ────────────────────────────────────────────────────────
\usepackage{graphicx}
\usepackage[dvipsnames,svgnames,table]{xcolor}
\usepackage{tikz}
\usetikzlibrary{arrows.meta, shapes.geometric, positioning,
                calc, backgrounds, patterns, decorations.pathmorphing,
                matrix}
\usepackage{pgfplots}
\pgfplotsset{compat=1.18}
\usepackage{tcolorbox}
\tcbuselibrary{skins, breakable, theorems}

% ─── Layout ───────────────────────────────────────────────────────────────────
\usepackage{titlesec}
\usepackage{fancyhdr}
\usepackage{enumitem}
\usepackage{booktabs}
\usepackage{array}
\usepackage{multicol}
\usepackage{marginnote}
\usepackage{mdframed}
\usepackage{caption}
\usepackage{subcaption}
\usepackage{parskip}
\usepackage{siunitx}
\usepackage[colorlinks=true, linkcolor=BrickRed,
            urlcolor=NavyBlue, citecolor=ForestGreen]{hyperref}

% ══════════════════════════════════════════════════════════════════════════════
% COLOUR PALETTE
% ══════════════════════════════════════════════════════════════════════════════
\definecolor{JourneyColor}{HTML}{065F46}   % forest green — journey headings
\definecolor{ModuleColor}{HTML}{0F766E}    % deep teal — module headings
\definecolor{PuzzleColor}{HTML}{B45309}    % deep amber — puzzle boxes
\definecolor{DiscoveryColor}{HTML}{D97706} % golden amber — discovery boxes
\definecolor{WarningColor}{HTML}{E11D48}   % coral red — "model breaks" boxes
\definecolor{ProblemColor}{HTML}{0F766E}   % dark teal — problem sets
\definecolor{ReflectColor}{HTML}{047857}   % green — reflection boxes
\definecolor{HistoryColor}{HTML}{9A3412}   % brown — history notes
\definecolor{LightBg}{HTML}{F8FAFC}        % near-white background

% ══════════════════════════════════════════════════════════════════════════════
% TCOLORBOX ENVIRONMENTS
% ══════════════════════════════════════════════════════════════════════════════
\newtcolorbox{puzzlebox}[1][]{
  enhanced, breakable,
  colback=PuzzleColor!6, colframe=PuzzleColor,
  fonttitle=\bfseries\sffamily, title={The Puzzle},
  attach boxed title to top left={yshift=-2mm, xshift=6mm},
  boxed title style={colback=PuzzleColor, colframe=PuzzleColor},
  borderline west={3pt}{0pt}{PuzzleColor}, #1}

\newtcolorbox{discoverybox}[1][]{
  enhanced, breakable,
  colback=DiscoveryColor!8, colframe=DiscoveryColor,
  fonttitle=\bfseries\sffamily, title={Discovery},
  attach boxed title to top left={yshift=-2mm, xshift=6mm},
  boxed title style={colback=DiscoveryColor, colframe=DiscoveryColor},
  borderline west={3pt}{0pt}{DiscoveryColor}, #1}

\newtcolorbox{problembox}[1][]{
  enhanced, breakable,
  colback=ProblemColor!5, colframe=ProblemColor,
  fonttitle=\bfseries\sffamily, title={Problems},
  attach boxed title to top left={yshift=-2mm, xshift=6mm},
  boxed title style={colback=ProblemColor, colframe=ProblemColor},
  borderline west={3pt}{0pt}{ProblemColor}, #1}

\newtcolorbox{deepbox}[1][]{
  enhanced, breakable,
  colback=JourneyColor!5, colframe=JourneyColor,
  fonttitle=\bfseries\sffamily, title={Deep Problems},
  attach boxed title to top left={yshift=-2mm, xshift=6mm},
  boxed title style={colback=JourneyColor, colframe=JourneyColor},
  borderline west={4pt}{0pt}{JourneyColor}, #1}

\newtcolorbox{historybox}[1][]{
  enhanced, breakable,
  colback=HistoryColor!7, colframe=HistoryColor!60,
  fonttitle=\bfseries\sffamily\small, title={Historical Note},
  attach boxed title to top left={yshift=-2mm, xshift=6mm},
  boxed title style={colback=HistoryColor!80, colframe=HistoryColor}, #1}

\newtcolorbox{reflectionbox}[1][]{
  enhanced, breakable,
  colback=ReflectColor!6, colframe=ReflectColor,
  fonttitle=\bfseries\sffamily, title={Reflection},
  attach boxed title to top left={yshift=-2mm, xshift=6mm},
  boxed title style={colback=ReflectColor, colframe=ReflectColor},
  borderline west={3pt}{0pt}{ReflectColor}, #1}

\newtcolorbox{modelbreaks}[1][]{
  enhanced, breakable,
  colback=WarningColor!6, colframe=WarningColor,
  fonttitle=\bfseries\sffamily, title={Where This Model Breaks},
  attach boxed title to top left={yshift=-2mm, xshift=6mm},
  boxed title style={colback=WarningColor, colframe=WarningColor},
  borderline west={4pt}{0pt}{WarningColor}, #1}

\newtcolorbox{namedconcept}[2][]{
  enhanced,
  colback=ModuleColor!6, colframe=ModuleColor,
  fonttitle=\bfseries\sffamily\large, title={#2},
  attach boxed title to top left={yshift=-2mm, xshift=6mm},
  boxed title style={colback=ModuleColor, colframe=ModuleColor}, #1}

% ══════════════════════════════════════════════════════════════════════════════
% THEOREM ENVIRONMENTS
% ══════════════════════════════════════════════════════════════════════════════
\theoremstyle{definition}
\newtheorem{observation}{Observation}[chapter]
\newtheorem{question}{Question}[chapter]
\newtheorem{experiment}{Experiment}[chapter]

% ══════════════════════════════════════════════════════════════════════════════
% SECTION FORMATTING
% ══════════════════════════════════════════════════════════════════════════════
\titleformat{\chapter}[display]
  {\normalfont\sffamily\huge\bfseries\color{JourneyColor}}
  {\textcolor{JourneyColor!60}{\chaptertitlename\ \thechapter}}
  {10pt}{\Huge}[\vspace{2pt}\textcolor{JourneyColor}{\titlerule[1.5pt]}]

\titleformat{\section}
  {\normalfont\sffamily\Large\bfseries\color{ModuleColor}}
  {\thesection}{1em}{}[\textcolor{ModuleColor!40}{\titlerule[0.6pt]}]

\titleformat{\subsection}
  {\normalfont\sffamily\large\bfseries\color{JourneyColor!80}}
  {\thesubsection}{1em}{}

\titleformat{\subsubsection}
  {\normalfont\sffamily\normalsize\bfseries}
  {\thesubsubsection}{1em}{}

% ══════════════════════════════════════════════════════════════════════════════
% HEADER / FOOTER
% ══════════════════════════════════════════════════════════════════════════════
\pagestyle{fancy}
\fancyhf{}
\fancyhead[LE]{\small\sffamily\textcolor{JourneyColor}{\leftmark}}
\fancyhead[RO]{\small\sffamily\textcolor{ModuleColor}{\rightmark}}
\fancyfoot[C]{\small\thepage}
\renewcommand{\headrulewidth}{0.4pt}

% ══════════════════════════════════════════════════════════════════════════════
% CUSTOM COMMANDS
% ══════════════════════════════════════════════════════════════════════════════
\newcommand{\chem}[1]{\ce{#1}}
\newcommand{\important}[1]{\textcolor{ModuleColor}{\textbf{#1}}}
\newcommand{\newterm}[1]{\textit{\textcolor{DiscoveryColor}{#1}}}
\newcommand{\sidenote}[1]{%
  \marginnote{\footnotesize\sffamily\textcolor{JourneyColor!70}{#1}}}

% ══════════════════════════════════════════════════════════════════════════════
% DOCUMENT
% ══════════════════════════════════════════════════════════════════════════════
\begin{document}
\setcounter{chapter}{0}
\chapter[Journey 1 --- What Is Stuff?]
        {Journey One\\[8pt]
         \textcolor{ModuleColor}{\large What Is Stuff?\ \ Matter and Its Patterns}}
\setcounter{section}{3}

% ══════════════════════════════════════════════════════════════════════════════
%  MODULE 1.4 — GASES, VOLUMES, AND AVOGADRO'S BRILLIANT HYPOTHESIS
% ══════════════════════════════════════════════════════════════════════════════
\section{Module 1.4 --- Gases, Volumes, and Avogadro's Brilliant Hypothesis}
\label{sec:mod1.4}

\subsection*{Motivation}

At the end of Module~1.3, Dalton had the right idea
but the wrong details.
He assumed water was \ce{HO} and got the atomic mass of oxygen wrong.
He had no way to determine whether a compound was \ce{HO} or \ce{H2O}
from mass ratios alone --- because either formula is consistent
with a 1:8 mass ratio if you choose the right atomic masses.

\begin{center}
\begin{tcolorbox}[
  enhanced, colback=LightBg, colframe=JourneyColor!40,
  width=0.86\textwidth, arc=3pt, fontupper=\sffamily
]
\centering
\textbf{The ambiguity:}\\[6pt]
\small
If water is \ce{HO}, then oxygen is 8 times heavier than hydrogen.\\
If water is \ce{H2O}, then oxygen is 16 times heavier than hydrogen.\\[6pt]
\normalsize
Both are consistent with the measured mass ratio H:O = 1:8.\\[4pt]
The mass data alone cannot distinguish between them.\\
We need a completely \emph{different kind of measurement}.
\end{tcolorbox}
\end{center}

\vspace{0.5em}

That different measurement came from gases.
Gases have a remarkable property:
their \emph{volumes}, at fixed temperature and pressure,
combine in simple ratios --- just as their masses do.
And those volume ratios turn out to be the key that unlocks
the true molecular formulas.

% ─── Section: The Behaviour of Gases ─────────────────────────────────────────
\subsection{The Behaviour of Ideal Gases}

Before we reach molecular formulas, we need to understand
how gases behave --- because the volume measurements
that unlock atomic theory all involve gases.

\subsubsection{Boyle's Law: Pressure and Volume}

Robert Boyle (1662) performed a simple but decisive experiment.
He trapped a fixed amount of gas in a tube
and systematically changed the pressure by adding mercury.

\begin{center}
\begin{tikzpicture}[
  every node/.style={font=\sffamily\small},
  scale=0.92
]
  \begin{axis}[
    width=7cm, height=6cm,
    xlabel={Volume $V$ (arbitrary units)},
    ylabel={Pressure $P$ (arbitrary units)},
    xmin=0.5, xmax=5.5,
    ymin=0, ymax=12,
    xtick={1,2,3,4,5},
    ytick={2,4,6,8,10},
    grid=both,
    grid style={line width=0.3pt, draw=JourneyColor!20},
    axis line style={color=JourneyColor!60},
    tick style={color=JourneyColor!60},
    title style={font=\sffamily\small\bfseries, color=JourneyColor},
    title={Boyle's Law: $P \propto 1/V$},
    label style={font=\sffamily\footnotesize}
  ]
    \addplot[
      thick, color=DiscoveryColor,
      domain=0.6:5.5, samples=80
    ]{6/x};
    \addplot[
      only marks, mark=*, mark size=2pt,
      color=JourneyColor
    ] coordinates {
      (1,6)(1.5,4)(2,3)(3,2)(6,1)
    };
  \end{axis}
  \node[font=\sffamily\footnotesize\itshape, text=HistoryColor]
    at (5.5, 0.3) {$PV = \text{constant}$};
\end{tikzpicture}
\hspace{1cm}
\begin{tikzpicture}[
  every node/.style={font=\sffamily\small},
  scale=0.92
]
  \begin{axis}[
    width=7cm, height=6cm,
    xlabel={Temperature $T$ (K)},
    ylabel={Volume $V$ (arbitrary units)},
    xmin=0, xmax=600,
    ymin=0, ymax=6,
    xtick={100,200,300,400,500},
    grid=both,
    grid style={line width=0.3pt, draw=JourneyColor!20},
    axis line style={color=JourneyColor!60},
    tick style={color=JourneyColor!60},
    title style={font=\sffamily\small\bfseries, color=JourneyColor},
    title={Charles's Law: $V \propto T$},
    label style={font=\sffamily\footnotesize}
  ]
    \addplot[
      thick, color=ModuleColor,
      domain=50:580, samples=80
    ]{x/100};
    \addplot[
      only marks, mark=*, mark size=2pt,
      color=JourneyColor
    ] coordinates {
      (100,1)(200,2)(273,2.73)(300,3)(400,4)(500,5)
    };
  \end{axis}
  \node[font=\sffamily\footnotesize\itshape, text=HistoryColor]
    at (4.5, 0.3) {$V/T = \text{constant}$};
\end{tikzpicture}
\end{center}

\vspace{0.5em}

Combining Boyle's and Charles's laws, and adding Avogadro's
discovery (which we are about to reach), produces the
\newterm{ideal gas law}:

\begin{center}
\begin{tcolorbox}[
  enhanced, colback=ModuleColor!6, colframe=ModuleColor,
  width=0.6\textwidth, arc=3pt, halign=center
]
\Large
$PV = nRT$
\end{tcolorbox}
\end{center}

where $P$ is pressure, $V$ is volume, $n$ is moles of gas,
$T$ is temperature in kelvin, and $R$ is the universal gas constant
($R = \SI{8.314}{\joule\per\mol\per\kelvin}$).

The crucial thing to notice: at fixed $T$ and $P$,
the volume depends \emph{only} on the number of moles ---
not on which gas it is.

\begin{center}
\large $V = \dfrac{nRT}{P}$\quad
$\Rightarrow$\quad equal moles of any gas occupy equal volumes.
\end{center}

\begin{problembox}
\begin{enumerate}[label=\textbf{\arabic*.}]
  \item A balloon contains \SI{2.0}{\mol} of helium at
        \SI{25}{\degreeCelsius} and \SI{1.0}{\text{atm}}.
        \begin{enumerate}[label=(\alph*)]
          \item Convert \SI{25}{\degreeCelsius} to kelvin.
          \item Calculate the volume of the balloon in litres.
                (Use $R = \SI{0.0821}{L\cdot atm\cdot mol^{-1}\cdot K^{-1}}$.)
          \item If the temperature drops to \SI{0}{\degreeCelsius}
                (at constant pressure), what is the new volume?
        \end{enumerate}

  \item A gas occupies \SI{4.0}{\litre} at \SI{2.0}{\text{atm}}
        and \SI{300}{\kelvin}.
        \begin{enumerate}[label=(\alph*)]
          \item How many moles of gas are present?
          \item What pressure would compress this gas
                to \SI{1.0}{\litre} at constant temperature?
          \item What temperature would double the volume
                at constant pressure?
        \end{enumerate}

  \item At standard temperature and pressure (STP:
        \SI{0}{\degreeCelsius}, \SI{1}{\text{atm}}),
        one mole of any ideal gas occupies \SI{22.4}{\litre}.
        \begin{enumerate}[label=(\alph*)]
          \item Verify this using the ideal gas law.
          \item A gas sample at STP occupies \SI{5.60}{\litre}.
                How many moles is this?
                How many molecules?
        \end{enumerate}
\end{enumerate}
\end{problembox}

% ─── Section: Gay-Lussac's Law of Combining Volumes ──────────────────────────
\subsection{Gay-Lussac's Discovery: Volumes React in Whole Numbers}

In 1809, Joseph Louis Gay-Lussac made an observation
that seems simple but carries enormous implications.

He was studying reactions between gases and noticed that
the volumes of reacting gases --- measured at the same
temperature and pressure --- were always in simple
whole-number ratios.

\begin{center}
\begin{tabular}{llccc}
\toprule
\textbf{Reaction} & \textbf{Equation} &
\textbf{Vol.\ gas 1} & \textbf{Vol.\ gas 2} &
\textbf{Vol.\ product} \\
\midrule
Hydrogen + oxygen & \ce{H2 + O2 -> H2O} &
2 parts \ce{H2} & 1 part \ce{O2} & 2 parts steam \\
Hydrogen + chlorine & \ce{H2 + Cl2 -> HCl} &
1 part \ce{H2} & 1 part \ce{Cl2} & 2 parts \ce{HCl} \\
Nitrogen + hydrogen & \ce{N2 + H2 -> NH3} &
1 part \ce{N2} & 3 parts \ce{H2} & 2 parts \ce{NH3} \\
Carbon monoxide + oxygen & \ce{CO + O2 -> CO2} &
2 parts \ce{CO} & 1 part \ce{O2} & 2 parts \ce{CO2} \\
\bottomrule
\end{tabular}
\end{center}

\vspace{0.5em}

This is the law of combining volumes.
And it poses an immediate puzzle.

\begin{puzzlebox}
Look at the reaction of hydrogen and chlorine:
\begin{center}
1 volume \ce{H2} + 1 volume \ce{Cl2} $\longrightarrow$ 2 volumes \ce{HCl}
\end{center}

Dalton's atomic model says atoms are indivisible.
Suppose the hydrogen and chlorine particles are single atoms
$h$ and $c$.
Then:
\begin{center}
1 particle $h$ + 1 particle $c$ $\longrightarrow$ 1 particle $hc$
\end{center}

If one particle of each produces one particle of product,
the volume should also go from 2 parts to 1 part.
But the experiment shows 2 parts going to \emph{2 parts}.

\begin{enumerate}[label=(\alph*), leftmargin=*]
  \item What is wrong with assuming hydrogen and chlorine
        exist as single atoms?

  \item What if hydrogen and chlorine are \emph{diatomic}
        molecules (\ce{H2} and \ce{Cl2})?
        Draw out the particle picture for:
        \begin{center}
        \ce{H2} + \ce{Cl2} $\longrightarrow$ 2 \ce{HCl}
        \end{center}
        How many particles go in, and how many come out?
        Does this explain the 1:1:2 volume ratio?

  \item Dalton's original theory said atoms are indivisible
        and cannot be split.
        But in the picture above, the \ce{H2} molecule
        \emph{is} split: each hydrogen atom goes to a different
        \ce{HCl} molecule.
        Is this a contradiction?
        What is the difference between splitting a \emph{molecule}
        and splitting an \emph{atom}?
\end{enumerate}
\end{puzzlebox}

\vspace{0.5em}

This puzzle was not fully resolved in Gay-Lussac's time.
Dalton himself rejected Gay-Lussac's data because it
contradicted his assumption of indivisible single-atom gases.
The resolution came from an Italian scientist
who made one of the most underappreciated contributions
in the history of chemistry.

% ─── Section: Avogadro's Hypothesis ──────────────────────────────────────────
\subsection{Avogadro's Hypothesis: Equal Volumes, Equal Numbers}

In 1811, Amedeo Avogadro proposed a hypothesis so elegant
it resolves all the contradictions at once.

\begin{namedconcept}{Avogadro's Hypothesis (1811)}
Equal volumes of all gases, at the same temperature
and pressure, contain equal \emph{numbers of molecules}.
\end{namedconcept}

\vspace{0.5em}

This hypothesis has two radical implications hidden inside it:

\textbf{First:} Gases like hydrogen, oxygen, and nitrogen
exist not as single atoms but as \emph{diatomic molecules}
--- \ce{H2}, \ce{O2}, \ce{N2}.
This is necessary to make Gay-Lussac's volume ratios work.

\textbf{Second:} The formula of water is \ce{H2O}, not \ce{HO}.
This finally resolves the ambiguity that trapped Dalton.

Let us see how the water formula falls out:

\begin{center}
\begin{tikzpicture}[
  every node/.style={font=\sffamily\small, align=center},
  atom/.style={circle, draw, thick, minimum size=0.6cm,
               font=\sffamily\tiny\bfseries},
  mol/.style={draw=JourneyColor!40, fill=JourneyColor!5,
              rounded corners=4pt, minimum height=1.2cm,
              inner sep=4pt}
]
  % Title row
  \node[font=\sffamily\small\bfseries, color=JourneyColor]
    at (0, 3.2) {Reactants};
  \node[font=\sffamily\small\bfseries, color=JourneyColor]
    at (8.5, 3.2) {Products};

  %--- Row 1: 2 H2 molecules ---
  \node[mol, minimum width=1.6cm] (hh1) at (-2.5, 1.8) {
    \tikz{
      \node[atom, fill=blue!20, draw=blue!70] (a) at (0,0) {H};
      \node[atom, fill=blue!20, draw=blue!70] (b) at (0.7,0) {H};
      \draw[thick, blue!50] (a)--(b);
    }
  };
  \node[mol, minimum width=1.6cm] (hh2) at (0.5, 1.8) {
    \tikz{
      \node[atom, fill=blue!20, draw=blue!70] (a) at (0,0) {H};
      \node[atom, fill=blue!20, draw=blue!70] (b) at (0.7,0) {H};
      \draw[thick, blue!50] (a)--(b);
    }
  };
  \node[font=\sffamily\footnotesize, text=JourneyColor!70]
    at (-1.0, 0.9) {2 volumes \ce{H2}};

  %--- Plus sign ---
  \node[font=\Large] at (2.2, 1.8) {$+$};

  %--- 1 O2 molecule ---
  \node[mol, minimum width=1.6cm] (oo) at (3.9, 1.8) {
    \tikz{
      \node[atom, fill=red!20, draw=red!70] (a) at (0,0) {O};
      \node[atom, fill=red!20, draw=red!70] (b) at (0.7,0) {O};
      \draw[thick, red!50] (a)--(b);
      \draw[thick, red!50] ([yshift=2pt]a.center)--([yshift=2pt]b.center);
    }
  };
  \node[font=\sffamily\footnotesize, text=JourneyColor!70]
    at (3.9, 0.9) {1 volume \ce{O2}};

  %--- Arrow ---
  \draw[-{Stealth[length=7pt]}, ultra thick, color=DiscoveryColor]
    (5.6, 1.8) -- (6.6, 1.8) node[above, midway, font=\sffamily\tiny]{react};

  %--- 2 H2O molecules ---
  \node[mol, minimum width=1.8cm] (w1) at (7.9, 1.8) {
    \tikz{
      \node[atom, fill=red!20, draw=red!70] (O) at (0.35,0.3) {O};
      \node[atom, fill=blue!20, draw=blue!70] (H1) at (0,  -0.3) {H};
      \node[atom, fill=blue!20, draw=blue!70] (H2) at (0.7,-0.3) {H};
      \draw[thick, purple!60] (O)--(H1);
      \draw[thick, purple!60] (O)--(H2);
    }
  };
  \node[mol, minimum width=1.8cm] (w2) at (10.2, 1.8) {
    \tikz{
      \node[atom, fill=red!20, draw=red!70] (O) at (0.35,0.3) {O};
      \node[atom, fill=blue!20, draw=blue!70] (H1) at (0,  -0.3) {H};
      \node[atom, fill=blue!20, draw=blue!70] (H2) at (0.7,-0.3) {H};
      \draw[thick, purple!60] (O)--(H1);
      \draw[thick, purple!60] (O)--(H2);
    }
  };
  \node[font=\sffamily\footnotesize, text=JourneyColor!70]
    at (9.05, 0.9) {2 volumes \ce{H2O} (steam)};

  % Volume count annotation
  \node[font=\sffamily\footnotesize\bfseries, text=HistoryColor]
    at (4.0, -0.1)
    {Volumes: 2 + 1 = 3 volumes in $\longrightarrow$ 2 volumes out};
  \node[font=\sffamily\footnotesize\itshape, text=ModuleColor]
    at (4.0, -0.7)
    {Each \ce{H2} contributes both H atoms to \emph{different} water molecules.
     Therefore \ce{H2} must split.};
\end{tikzpicture}
\end{center}

\vspace{0.5em}

This diagram reveals why the water formula \emph{must} be \ce{H2O}
and not \ce{HO}.
If water were \ce{HO}, one volume of hydrogen plus one volume
of oxygen would give exactly one volume of water --- but
experiments consistently show two volumes of hydrogen
plus one volume of oxygen give \emph{two} volumes of steam.
The factor of two in the product demands two hydrogen atoms per water.

\begin{discoverybox}
\textbf{What Avogadro's hypothesis fixes:}
\begin{itemize}[nosep, leftmargin=2em]
  \item[\checkmark] Water is \ce{H2O}, not \ce{HO}
  \item[\checkmark] Oxygen atomic mass is 16, not 8
  \item[\checkmark] Hydrogen, oxygen, nitrogen, chlorine are diatomic
        (\ce{H2}, \ce{O2}, \ce{N2}, \ce{Cl2})
  \item[\checkmark] Volume ratios in gas reactions equal mole ratios
  \item[\checkmark] Connects microscopic (molecule count) to
        macroscopic (volume) measurements
\end{itemize}
\end{discoverybox}

\begin{historybox}
\textbf{Why Avogadro Was Ignored for 50 Years}

Avogadro published his hypothesis in 1811.
It was almost completely ignored until 1860,
when Stanislao Cannizzaro presented a systematic argument
for it at the First International Chemical Congress in Karlsruhe.

Why ignored?
\begin{itemize}[nosep, leftmargin=2em]
  \item Dalton actively opposed it (it contradicted his model).
  \item The idea of diatomic molecules (\ce{H2}, \ce{O2}) seemed
        strange --- why would two identical atoms stick together?
  \item There was no theoretical framework to explain molecular bonding.
  \item Chemistry in 1811 lacked international communication
        and consensus-building mechanisms.
\end{itemize}

The Karlsruhe Congress of 1860 is one of the most important
events in the history of chemistry.
Within a decade, Cannizzaro's arguments produced consistent
atomic masses, which in turn allowed Mendeleev to order
the elements by atomic mass and discover the periodic table.

The lesson: even a correct idea can take half a century
to gain acceptance if it lacks the right theoretical
context and the right champions.
\end{historybox}

\begin{problembox}
\begin{enumerate}[label=\textbf{\arabic*.}]
  \item Use Avogadro's hypothesis and Gay-Lussac's volume ratios
        to determine the formula of ammonia.
        The experiment shows:
        \begin{center}
        1 volume \ce{N2} + 3 volumes \ce{H2} $\rightarrow$
        2 volumes ammonia
        \end{center}
        \begin{enumerate}[label=(\alph*)]
          \item If both \ce{N2} and \ce{H2} are diatomic,
                how many nitrogen atoms and hydrogen atoms
                are in 2 molecules of ammonia?
          \item Therefore, how many nitrogen and hydrogen
                atoms are in \emph{one} molecule of ammonia?
          \item What is the molecular formula of ammonia?
          \item Verify: does this match the known formula \ce{NH3}?
        \end{enumerate}

  \item Use the same reasoning for hydrogen chloride.
        The experiment shows:
        \begin{center}
        1 volume \ce{H2} + 1 volume \ce{Cl2} $\rightarrow$
        2 volumes \ce{HCl}
        \end{center}
        \begin{enumerate}[label=(\alph*)]
          \item Why does 1 volume + 1 volume give 2 volumes
                (not 1 volume)?
          \item Draw the particle picture showing how
                \ce{H2} and \ce{Cl2} split and recombine.
          \item What does this tell us about the behaviour
                of diatomic molecules in reactions?
        \end{enumerate}

  \item At \SI{0}{\degreeCelsius} and \SI{1}{\text{atm}}:
        \begin{center}
        \SI{2.0}{\litre} \ce{H2} + \SI{1.0}{\litre} \ce{O2}
        $\rightarrow$ \SI{2.0}{\litre} steam
        \end{center}
        \begin{enumerate}[label=(\alph*)]
          \item How many moles of \ce{H2} are present?
          \item How many moles of \ce{O2}?
          \item How many moles of steam are produced?
          \item What is the mass of steam produced?
        \end{enumerate}
\end{enumerate}
\end{problembox}

% ─── Section: Stoichiometry ───────────────────────────────────────────────────
\subsection{Stoichiometry: The Arithmetic of Reactions}

With Avogadro's hypothesis in hand, we can now do
something powerful: use a balanced chemical equation
as a quantitative map of a reaction.

The coefficients in a balanced equation tell us
the \emph{mole ratios} in which substances react and are produced.

\begin{center}
\begin{tcolorbox}[
  enhanced, colback=JourneyColor!6, colframe=JourneyColor,
  width=0.88\textwidth, arc=3pt
]
\centering\sffamily\small
\ce{2H2 + O2 -> 2H2O}\\[6pt]
\normalsize
2 moles \ce{H2} + 1 mole \ce{O2} $\longrightarrow$ 2 moles \ce{H2O}\\
\small
(or 2 volumes + 1 volume $\longrightarrow$ 2 volumes, for gases at equal $T$, $P$)\\[4pt]
\normalsize
\SI{4.0}{\gram} \ce{H2} + \SI{32.0}{\gram} \ce{O2}
$\longrightarrow$ \SI{36.0}{\gram} \ce{H2O}
\end{tcolorbox}
\end{center}

\vspace{0.5em}

This map connects moles to masses to volumes.
The systematic procedure for using it is called
\newterm{stoichiometry} (from the Greek for
``element measure'').

\subsubsection{The Stoichiometry Procedure}

\begin{center}
\begin{tikzpicture}[
  every node/.style={font=\sffamily\small, align=center},
  box/.style={draw=JourneyColor!50, fill=JourneyColor!8,
              rounded corners=3pt, minimum width=2.4cm,
              minimum height=0.9cm},
  arr/.style={-{Stealth[length=5pt]}, thick, color=ModuleColor}
]
  \node[box] (given) at (0,0)    {Given\\(mass or volume)};
  \node[box] (mol1)  at (3.5,0)  {Convert to\\moles};
  \node[box] (mol2)  at (7.0,0)  {Use mole\\ratio};
  \node[box] (final) at (10.5,0) {Convert to\\desired units};

  \draw[arr] (given) -- (mol1)
    node[midway, above, font=\sffamily\footnotesize]{$\div M$};
  \draw[arr] (mol1)  -- (mol2)
    node[midway, above, font=\sffamily\footnotesize]{coefficients};
  \draw[arr] (mol2)  -- (final)
    node[midway, above, font=\sffamily\footnotesize]{$\times M$\\or $\times 22.4$};
\end{tikzpicture}
\end{center}

\vspace{0.5em}

\subsubsection{Worked Example: Haber Process}

The Haber process synthesises ammonia from nitrogen and hydrogen:
\begin{center}
\ce{N2 + 3H2 -> 2NH3}
\end{center}

\textbf{Question:} What mass of ammonia is produced from
\SI{28.0}{\gram} of nitrogen, assuming excess hydrogen?

\textbf{Step 1: Moles of \ce{N2}}
\[
n(\ce{N2}) = \frac{\SI{28.0}{\gram}}{\SI{28.0}{\gram\per\mol}} = \SI{1.00}{\mol}
\]

\textbf{Step 2: Mole ratio}

From the equation: 1 mol \ce{N2} produces 2 mol \ce{NH3}.
\[
n(\ce{NH3}) = 1.00 \times 2 = \SI{2.00}{\mol}
\]

\textbf{Step 3: Mass of \ce{NH3}}
\[
m(\ce{NH3}) = \SI{2.00}{\mol} \times \SI{17.0}{\gram\per\mol} = \SI{34.0}{\gram}
\]

\subsubsection{Limiting Reagents}

In practice, reactants are rarely present in exact stoichiometric
amounts.
The \newterm{limiting reagent} is the reactant that runs out first
and determines how much product forms.

\textbf{Example:} \SI{10.0}{\gram} \ce{N2} reacts with
\SI{10.0}{\gram} \ce{H2}.
Which is limiting?

\ce{N2}: $n = 10.0/28.0 = \SI{0.357}{\mol}$\\
\ce{H2}: $n = 10.0/2.0  = \SI{5.00}{\mol}$

The equation requires 3 mol \ce{H2} per mol \ce{N2}.
For 0.357 mol \ce{N2}, we need $0.357 \times 3 = \SI{1.07}{\mol}$ \ce{H2}.
We have \SI{5.00}{\mol} \ce{H2} --- more than enough.
Therefore \ce{N2} is the limiting reagent.

Product: $0.357 \times 2 = \SI{0.714}{\mol}\ \ce{NH3}
= \SI{12.1}{\gram}$.

\begin{problembox}
\begin{enumerate}[label=\textbf{\arabic*.}]
  \item Iron(III) oxide reacts with carbon monoxide
        in the blast furnace:
        \begin{center}
        \ce{Fe2O3 + 3CO -> 2Fe + 3CO2}
        \end{center}
        \begin{enumerate}[label=(\alph*)]
          \item What mass of iron is produced from
                \SI{160}{\gram} of \ce{Fe2O3}?
          \item What volume of \ce{CO2} (at STP) is produced
                alongside the iron in part (a)?
          \item If only \SI{88}{\gram} of iron is actually
                recovered in practice, what is the
                \newterm{percentage yield}?
        \end{enumerate}

  \item Silicon is produced by the reaction:
        \begin{center}
        \ce{SiO2 + 2C -> Si + 2CO}
        \end{center}
        \SI{60.1}{\gram} of \ce{SiO2} reacts with
        \SI{36.0}{\gram} of carbon.
        \begin{enumerate}[label=(\alph*)]
          \item Calculate moles of each reactant.
          \item Determine the limiting reagent.
          \item Calculate the theoretical yield of silicon in grams.
          \item The actual yield is \SI{22.4}{\gram}.
                What is the percentage yield?
        \end{enumerate}

  \item An unknown gas \textbf{X} reacts with \ce{O2}:
        \begin{center}
        \ce{X + O2 -> XO2}
        \end{center}
        \SI{32.0}{\gram} of \textbf{X} reacts with exactly
        \SI{32.0}{\gram} of \ce{O2}.
        \begin{enumerate}[label=(\alph*)]
          \item From the stoichiometry, how many moles of
                \textbf{X} must react with each mole of \ce{O2}?
          \item What is the molar mass of \textbf{X}?
          \item Identify element \textbf{X}.
        \end{enumerate}
\end{enumerate}
\end{problembox}

% ─── Section: Connecting the Laws ────────────────────────────────────────────
\subsection{The Complete Picture: How the Laws Connect}

We can now draw the full map of Journey~1's discoveries.

\begin{center}
\begin{tikzpicture}[
  every node/.style={font=\sffamily\footnotesize, align=center},
  lawbox/.style={draw=JourneyColor!60, fill=JourneyColor!8,
                 rounded corners=4pt, minimum width=3.2cm,
                 minimum height=0.9cm},
  concbox/.style={draw=ModuleColor!70, fill=ModuleColor!8,
                  rounded corners=4pt, minimum width=3.2cm,
                  minimum height=0.9cm},
  predbox/.style={draw=DiscoveryColor!80, fill=DiscoveryColor!8,
                  rounded corners=4pt, minimum width=3.2cm,
                  minimum height=0.9cm},
  arr/.style={-{Stealth[length=4pt]}, thick, color=JourneyColor!60}
]
  % Empirical Laws column
  \node[font=\sffamily\small\bfseries, color=JourneyColor]
        at (0, 4.2) {Empirical Laws};
  \node[lawbox] (cm)  at (0, 3.2) {Conservation of mass\\(Lavoisier, 1785)};
  \node[lawbox] (dp)  at (0, 2.0) {Definite proportions\\(Proust, 1799)};
  \node[lawbox] (mp)  at (0, 0.8) {Multiple proportions\\(Dalton, 1803)};
  \node[lawbox] (cv)  at (0,-0.4) {Combining volumes\\(Gay-Lussac, 1809)};

  % Core Theory
  \node[font=\sffamily\small\bfseries, color=ModuleColor]
        at (5.5, 4.2) {Core Theory};
  \node[concbox, minimum width=4cm, minimum height=2.2cm]
        (at) at (5.5, 1.5)
        {\textbf{Atomic-Molecular Theory}\\[4pt]
         Matter = discrete atoms\\
         Compounds = fixed atom ratios\\
         Gases = equal vol = equal moles\\
         Elements may be diatomic};

  % Predictions column
  \node[font=\sffamily\small\bfseries, color=DiscoveryColor!80]
        at (11, 4.2) {Consequences};
  \node[predbox] (f1) at (11, 3.2) {Molecular formulas\\determinable};
  \node[predbox] (f2) at (11, 2.0) {Atomic mass scale\\self-consistent};
  \node[predbox] (f3) at (11, 0.8) {Stoichiometry:\\quantitative prediction};
  \node[predbox] (f4) at (11,-0.4) {Periodic pattern\\discoverable};

  % Arrows in
  \foreach \src in {cm, dp, mp, cv}
    \draw[arr] (\src.east) -- (at.west);

  % Arrows out
  \foreach \tgt in {f1, f2, f3, f4}
    \draw[arr] (at.east) -- (\tgt.west);
\end{tikzpicture}
\end{center}

\vspace{0.5em}

The last entry in the ``Consequences'' column --- a periodic
pattern becomes discoverable --- is the thread that leads into
Journey~3 (the periodic table) and eventually all of
modern chemistry.

% ─── Section: Deep Problems ───────────────────────────────────────────────────
\subsection*{Deep Problems}

\begin{deepbox}
\begin{enumerate}[label=\textbf{D\arabic*.}]

  \item \textbf{The Cannizzaro Argument.}

        At the 1860 Karlsruhe Congress, Stanislao Cannizzaro
        argued as follows to establish a consistent atomic mass scale.

        Consider a series of gaseous compounds of hydrogen.
        For each compound, measure the molar mass (by gas density)
        and the percentage of hydrogen by mass.
        The mass of hydrogen per mole of compound is then
        (molar mass) $\times$ (fraction that is hydrogen).

        Here is the data for several hydrogen compounds:

        \begin{center}
        \begin{tabular}{lrr}
        \toprule
        Compound & Molar mass & H fraction \\
        \midrule
        \ce{HCl}   & \SI{36.5}{\gram\per\mol}  & 2.74\% \\
        \ce{H2O}   & \SI{18.0}{\gram\per\mol}  & 11.11\% \\
        \ce{NH3}   & \SI{17.0}{\gram\per\mol}  & 17.65\% \\
        \ce{CH4}   & \SI{16.0}{\gram\per\mol}  & 25.00\% \\
        \ce{H2}    & \SI{2.0}{\gram\per\mol}   & 100\% \\
        \bottomrule
        \end{tabular}
        \end{center}

        \begin{enumerate}[label=(\alph*)]
          \item Calculate the mass of hydrogen per mole
                of each compound.
          \item What is the smallest value you get?
          \item Cannizzaro argued this smallest value is
                the mass of \emph{one hydrogen atom}.
                Why is that a reasonable inference?
          \item Use your answer to find the mass of a hydrogen atom.
          \item How does this fix the atomic mass scale
                that Dalton got wrong?
        \end{enumerate}

  \item \textbf{The Ideal Gas Law and Molecular Reality.}

        The ideal gas law $PV = nRT$ is derived from the assumption
        that gas molecules are point particles with no volume
        and no interactions between them.

        \begin{enumerate}[label=(\alph*)]
          \item Real gases deviate from the ideal gas law at
                high pressures and low temperatures.
                At \emph{high pressure}, the volume of real \ce{CO2}
                is \emph{larger} than $nRT/P$ predicts.
                What does this suggest about the actual volume
                of \ce{CO2} molecules?

          \item At \emph{low temperatures}, real gases are harder
                to compress than the ideal law predicts ---
                but then suddenly liquefy.
                What does this suggest about the forces between molecules?

          \item Van der Waals modified the ideal gas law to:
                \[
                  \left(P + \frac{an^2}{V^2}\right)(V - nb) = nRT
                \]
                The term $nb$ corrects for molecular volume.
                The term $an^2/V^2$ corrects for intermolecular
                attractions.
                For \ce{CO2}: $a = 3.59\,\text{L}^2\,\text{atm\,mol}^{-2}$,
                $b = 0.0427\,\text{L\,mol}^{-1}$.
                Calculate the pressure of \SI{1.00}{\mol} of \ce{CO2}
                at \SI{0}{\degreeCelsius} in a \SI{1.00}{\litre} vessel,
                using both the ideal law and van der Waals.
                Which gives a higher pressure, and why?
        \end{enumerate}

  \item \textbf{Density of a Gas and Its Molar Mass.}

        An unknown gas \textbf{Z} has a density of
        \SI{1.96}{\gram\per\litre} at \SI{0}{\degreeCelsius}
        and \SI{1.0}{\text{atm}}.

        \begin{enumerate}[label=(\alph*)]
          \item Using the fact that 1 mole of ideal gas
                occupies \SI{22.4}{\litre} at STP,
                find the molar mass of gas \textbf{Z}.

          \item The gas burns in oxygen to give only \ce{CO2}
                and \ce{H2O}, in the mass ratio
                \ce{CO2}:\ce{H2O} = 88:36 per \SI{44}{\gram}
                of gas burned.
                Find the empirical formula of \textbf{Z}.

          \item Does the empirical formula match the molar mass?
                What is the molecular formula?

          \item Identify \textbf{Z}.
                (Hint: it is a common fuel.)
        \end{enumerate}

\end{enumerate}
\end{deepbox}

% ─── Reflection ───────────────────────────────────────────────────────────────
\subsection*{Reflection}

\begin{reflectionbox}
\begin{enumerate}[label=\textbf{R\arabic*.}]
  \item Avogadro's hypothesis was published in 1811 and ignored
        until 1860.
        \textbf{What conditions are necessary for a correct
        scientific idea to be accepted?}
        Is correctness sufficient?
        What role do experimental confirmation,
        theoretical context, and personal authority play?

  \item The formula \ce{H2O} and the diatomic nature of \ce{H2}
        and \ce{O2} were established from \emph{volume} measurements,
        not mass measurements.
        \textbf{Why could mass measurements alone not settle this?}
        What additional type of evidence was needed, and why?

  \item Stoichiometry lets us predict the outcome of a reaction
        with precision --- but only if the reaction goes to completion
        and only one reaction pathway is possible.
        \textbf{In real industrial processes, neither assumption holds.}
        What additional concepts (beyond stoichiometry) would a
        chemical engineer need to actually design a production process?
        (Write a hypothesis --- we will develop this in Journey~5.)

  \item We built the entire Journey~1 picture from five
        empirical laws (conservation of mass, definite proportions,
        multiple proportions, combining volumes, Avogadro's hypothesis).
        \textbf{Are these laws \emph{derived} from the atomic theory,
        or does the atomic theory \emph{explain} them?}
        Is there a difference?
        Which direction of reasoning is more powerful for science?
\end{enumerate}
\end{reflectionbox}

\vspace{1.5em}
\begin{center}
\begin{tcolorbox}[
  enhanced, colback=DiscoveryColor!8, colframe=DiscoveryColor,
  width=0.88\textwidth, arc=3pt
]
\centering\sffamily\small
\textbf{What we invented in this module.}\\[4pt]
We invented the idea that \emph{volume is a proxy for particle count}.
This was Avogadro's genius: not that he proposed a new theory,
but that he found a way to count invisible particles
using a measurement anyone could make with a syringe and a ruler.
With this bridge --- volume to moles to atoms ---
the entire edifice of quantitative chemistry becomes accessible.
Every formula in a chemistry textbook, every industrial process calculation,
every pharmaceutical dosage, every materials synthesis protocol
ultimately traces back to the argument made in this module:
equal volumes of gases contain equal numbers of molecules,
and therefore stoichiometric ratios in equations
correspond to real particle ratios in nature.
\end{tcolorbox}
\end{center}

\vspace{2em}
\noindent\textcolor{JourneyColor!40}{\rule{\textwidth}{0.6pt}}
\begin{center}
\textit{\sffamily\small
Module 1.4 complete.\\
Journey~1 concludes with Module~1.5:
\textbf{Putting It Together --- From Observation to Atomic World}.\\
Or proceed directly to Journey~2:
\textbf{Architecture of the Atom}.}
\end{center}

\end{document}